\chapter{The Donation}

\lettrine{C}{olbert} reappeared beneath the curtains.

\zz
<Have you heard?> said Mazarin.

\zz
<Alas! yes, my lord.>

\zz
<Can he be right? Can all this money be badly acquired?>

\zz
<A Theatin, monseigneur, is a bad judge in matters of finance,>  replied Colbert, coolly. <And yet it is very possible that, according to  his theological views, your eminence has been, in a certain degree, in the  wrong. People generally find they have been so,—when they die.>

<In the first place, they commit the wrong of dying, Colbert.>

<That is true, my lord. Against whom, however, did the Theatin make out  that you had committed these wrongs? Against the king?>

Mazarin shrugged his shoulders. <As if I had not saved both his state and  his finances.>

<That admits of no contradiction, my lord.>

<Does it? Then I have received a merely legitimate salary, in spite of  the opinion of my confessor?>

<That is beyond doubt.>

<And I might fairly keep for my own family, which is so needy, a good  fortune,—the whole, even, of which I have earned?>

<I see no impediment to that, monseigneur.>

<I felt assured that in consulting you, Colbert, I should have good  advice,> replied Mazarin, greatly delighted.

Colbert resumed his pedantic look. <My lord,> interrupted he,  <I think it would be quite as well to examine whether what the Theatin  said is not a snare.>

<Oh! no; a snare? What for? The Theatin is an honest man.>

<He believed your eminence to be at death's door, because your  eminence consulted him. Did I not hear him say—<Distinguish that  which the king has given you from that which you have given yourself.>  Recollect, my lord, if he did not say something a little like that to  you?—that is quite a theatrical speech.>

<That is possible.>

<In which case, my lord, I should consider you as required by the Theatin  to\longdash>

<To make restitution!> cried Mazarin, with great warmth.

<Eh! I do not say no.>

<What, of all! You do not dream of such a thing! You speak just as the  confessor did.>

<To make restitution of a part,—that is to say, his majesty's  part; and that, monseigneur, may have its dangers. Your eminence is too  skilful a politician not to know that, at this moment, the king does not  possess a hundred and fifty thousand livres clear in his coffers.>

<That is not my affair,> said Mazarin, triumphantly; <that  belongs to M. le Surintendant Fouquet, whose accounts I gave you to verify some  months ago.>

Colbert bit his lips at the name of Fouquet. <His majesty,> said  he, between his teeth, <has no money but that which M. Fouquet collects:  your money, monseigneur, would afford him a delicious banquet.>

<Well, but I am not the superintendent of his majesty's  finances—I have my purse—surely I would do much for his  majesty's welfare—some legacy—but I cannot disappoint my  family.>

<The legacy of a part would dishonour you and offend the king. Leaving a  part to his majesty, is to avow that that part has inspired you with doubts as  to the lawfulness of the means of acquisition.>

<Monsieur Colbert!>

<I thought your eminence did me the honour to ask my advice?>

<Yes, but you are ignorant of the principal details of the  question.>

<I am ignorant of nothing, my lord; during ten years, all the columns of  figures which are found in France, have passed into review before me; and if I  have painfully nailed them into my brain, they are there now so well riveted,  that, from the office of M. Letellier, who is sober, to the little secret  largesses of M. Fouquet, who is prodigal, I could recite, figure by figure, all  the money that is spent in France from Marseilles to Cherbourg.>

<Then, you would have me throw all my money into the coffers of the  king!> cried Mazarin, ironically; and from whom, at the same time the  gout forced painful moans. <Surely the king would reproach me with  nothing, but he would laugh at me, while squandering my millions, and with good  reason.>

<Your eminence has misunderstood me. I did not, the least in the world,  pretend that his majesty ought to spend your money.>

<You said so, clearly, it seems to me, when you advised me to give it to  him.>

<Ah,> replied Colbert, <that is because your eminence,  absorbed as you are by your disease, entirely loses sight of the character of  Louis XIV>

<How so?>

<That character, if I may venture to express myself thus, resembles that  which my lord confessed just now to the Theatin.>

<Go on—that is?>

<Pride! Pardon me, my lord, haughtiness, nobleness; kings have no pride,  that is a human passion.>

<Pride,—yes, you are right. Next?>

<Well, my lord, if I have divined rightly, your eminence has but to give  all your money to the king, and that immediately.>

<But for what?> said Mazarin, quite bewildered.

<Because the king will not accept of the whole.>

<What, and he a young man, and devoured by ambition?>

<Just so.>

<A young man who is anxious for my death\longdash>

<My lord!>

<To inherit, yes, Colbert, yes; he is anxious for my death, in order to  inherit. Triple fool that I am! I would prevent him!>

<Exactly: if the donation were made in a certain form he would refuse  it.>

<Well; but how?>

<That is plain enough. A young man who has yet done nothing—who  burns to distinguish himself—who burns to reign alone, will never take  anything ready built, he will construct for himself. This prince, monseigneur,  will never be content with the Palais Royal, which M. de Richelieu left him,  nor with the Palais Mazarin, which you have had so superbly constructed, nor  with the Louvre, which his ancestors inhabited; nor with St Germain, where he  was born. All that does not proceed from himself, I predict, he will  disdain.>

<And you will guarantee, that if I give my forty millions to the  king\longdash>

<Saying certain things to him at the same time, I guarantee he will  refuse them.>

<But those things—what are they?>

<I will write them, if my lord will have the goodness to dictate  them.>

<Well, but, after all, what advantage will that be to me?>

<An enormous one. Nobody will afterwards be able to accuse your eminence  of that unjust avarice with which pamphleteers have reproached the most  brilliant mind of the present age.>

<You are right, Colbert, you are right; go, and seek the king, on my  part, and take him my will.>

<Your donation, my lord.>

<But, if he should accept it; if he should even think of accepting  it!>

<Then there would remain thirteen millions for your family, and that is a  good round sum.>

<But then you would be either a fool or a traitor.>

<And I am neither the one nor the other, my lord. You appear to be much  afraid that the king will accept; you have a deal more reason to fear that he  will not accept.>

<But, see you, if he does not accept, I should like to guarantee my  thirteen reserved millions to him—yes, I will do so—yes. But my  pains are returning, I shall faint. I am very, very ill, Colbert; I am near my  end!>

Colbert started. The cardinal was indeed very ill; large drops of sweat flowed  down upon his bed of agony, and the frightful pallor of a face streaming with  water was a spectacle which the most hardened practitioner could not have  beheld without much compassion. Colbert was, without doubt, very much affected,  for he quitted the chamber, calling Bernouin to attend to the dying man, and  went into the corridor. There, walking about with a meditative expression,  which almost gave nobility to his vulgar head, his shoulders thrown up, his  neck stretched out, his lips half open, to give vent to unconnected fragments  of incoherent thoughts, he lashed up his courage to the pitch of the  undertaking contemplated, whilst within ten paces of him, separated only by a  wall, his master was being stifled by anguish which drew from him lamentable  cries, thinking no more of the treasures of the earth, or of the joys of  Paradise, but much of all the horrors of hell. Whilst burning-hot napkins,  physic, revulsives, and Guenaud, who was recalled, were performing their  functions with increased activity, Colbert, holding his great head in both his  hands, to compress within it the fever of the projects engendered by the brain,  was meditating the tenor of the donation he would make Mazarin write, at the  first hour of respite his disease should afford him. It would appear as if all  the cries of the cardinal, and all the attacks of death upon this  representative of the past, were stimulants for the genius of this thinker with  the bushy eyebrows, who was turning already towards the rising sun of a  regenerated society. Colbert resumed his place at Mazarin's pillow at the  first interval of pain, and persuaded him to dictate a donation thus conceived.

\begin{mail}{}{}
About to appear before God, the Master of mankind, I beg the king, who  was my master on earth, to resume the wealth which his bounty has bestowed upon  me, and which my family would be happy to see pass into such illustrious hands.  The particulars of my property will be found—they are drawn up—at  the first requisition of his majesty, or at the last sigh of his most devoted  servant,
\closeletter{Jules, Cardinal de Mazarin.}
\end{mail}


The cardinal sighed heavily as he signed this; Colbert sealed the packet, and  carried it immediately to the Louvre, whither the king had returned.

He then went back to his own home, rubbing his hands with the confidence of  workman who has done a good day's work.  